\FloatBarrier
\section{Downstream beam tube optimization}

The remaining downstream contacts recorded in this historical study are resolved by the subsequent socket
revision in Section~\ref{sec:downstream-joints}.

The next beamline step covers the two downstream tube segments. Their engineering tessellations are already
coarse, so reductions are modest: 5,146 to 5,000 and 672 to 600 triangles. Together they decrease from 5,818
to 5,600 triangles (3.75\%) and from 291,068 to 280,168 bytes. Names, stainless-steel materials, coordinate
frames, and placements are retained in GEMC3 and GEMC2; no CAD database row is renamed or moved.

Segment 1 originally contains seven disconnected surface components, 392 boundary edges, fourteen boundary
loops, and 1,021 faces flagged by the mesh self-intersection filter. Generic hole closing caps the bore and
is unsuitable. The repair registers the seven inner/outer boundary-ring pairs on matching angular grids
while preserving their measured radii and axial positions, then adds 392 annular faces. All 392 seam
vertices are protected during decimation. The output is a single closed, manifold tube of genus one; the
passage remains open. Segment 2 retains its original single-component, genus-one topology. Neither output
has detected self-intersections or exact zero-area triangles.

Table~\ref{tab:downstream-meshes} summarizes the mesh inventory.

\begin{table}[htbp]
\centering
\caption{Downstream triangle counts and surface-component changes.}
\label{tab:downstream-meshes}
\begin{tabular}{lrrrr}
\toprule
Mesh & Original & Optimized & Original components & Output components \\
\midrule
Segment 1 & 5,146 & 5,000 & 7 & 1 \\
Segment 2 & 672 & 600 & 1 & 1 \\
\bottomrule
\end{tabular}
\end{table}

The largest bidirectionally sampled surface differences are 0.7454 mm and 0.4574 mm, mainly around short
connector and flange features. Bounds change by at most 0.011 mm. The 0.05 mm apex sampling limit does not
apply to these tubes. Every original/output vertex and approximately 200,000 edge/surface samples per
direction are checked; these sampled maxima are not proven global error bounds. Segment 1 has no
well-defined original volume because its surface is open. Decimation changes its annularly repaired
reference volume by -0.01632\%; this does not establish original-to-repaired mass equivalence. Segment 2
changes by -0.00170\%. The builder reproduces both outputs byte for byte from archived originals.

The pair-only benchmark uses the recorded GEMC3 executable and Geant4 stack, FTFP\_BERT without optical
processes, air, zero magnetic field, one thread, and identical default production cuts. GUI and event output
are disabled. Matched 5 GeV LUND primaries start at the target, uniform in solid angle from 0.1 to 1.3
degrees and over full azimuth. The smaller forward aperture illuminates the downstream tubes. Five
alternating-order repetitions use seeds 12345--12349 and 200 events per electron, proton, positive-pion, or
neutron run. Separate zero-event runs measure startup without overlap checking; elapsed wall time since
beamOn measures transport. CPU and peak RSS are retained separately. Absolute timings are not pooled with
the apex or LTCC workloads.

Table~\ref{tab:downstream-timing} gives the separate particle-transport measurements.

\begin{table}[htbp]
\centering
\caption{Median transport wall seconds per 200 events over five repetitions.}
\label{tab:downstream-timing}
\begin{tabular}{lrrr}
\toprule
Particle & Original & Optimized & Original/optimized \\
\midrule
e- & 13.008 & 11.401 & 1.14 \\
proton & 2.109 & 1.901 & 1.11 \\
pi+ & 2.318 & 2.161 & 1.07 \\
neutron & 1.340 & 1.204 & 1.11 \\
\bottomrule
\end{tabular}
\end{table}

Median startup simulation wall time changes from 0.6924 to 0.6925 s; loading peak RSS changes from 123.27 to
123.06 MiB. These small initialization differences should not be interpreted as a substantial startup or
memory improvement. The original/optimized median transport wall-time ratios range from 1.07 to 1.14. All 52
processes exit normally, and all forty transport runs report their 200 requested events. No event-abort
exception is recorded.

Table~\ref{tab:downstream-warnings} quantifies the retained geometry and navigation diagnostics.

\begin{table}[htbp]
\centering
\caption{Downstream construction and transport warning counts.}
\label{tab:downstream-warnings}
\begin{tabular}{lrr}
\toprule
Diagnostic & Original & Optimized \\
\midrule
Surface-hole warnings, five loading runs & 5 & 0 \\
e- stuck tracks, five runs & 10 & 14 \\
proton stuck tracks, five runs & 0 & 0 \\
pi+ stuck tracks, five runs & 0 & 4 \\
neutron stuck tracks, five runs & 0 & 1 \\
Placement reports, overlap sample & 1 & 3 \\
\bottomrule
\end{tabular}
\end{table}

The repair eliminates the original surface-hole import warning. Placement overlap reports and stuck-track
diagnostics remain, as quantified below. A completion footer and requested event count do not prove that
every track followed the intended geometry. The original mesh lacks closed surfaces, and the repaired model
changes the represented material boundaries. Therefore transport timing differences cannot be assigned
solely to triangle-count reduction or used as proof of faster correct transport. The unchanged engineering
placement overlap and the thin, coarsely tessellated connector surfaces require further geometry/navigation
work. Detector response and full-experiment performance are outside this study.

\FloatBarrier
Figure~\ref{fig:downstream-stl} compares the stored triangulations.

\begin{figure}[htbp]
\centering
\includegraphics[width=\textwidth]{figures/downstream_stl_comparison.pdf}
\caption{Original (left) and optimized (right) downstream tube segments in millimetre coordinates.
Both models retain the through-bore; the paired views use the same camera and scale.}
\label{fig:downstream-stl}
\end{figure}
